% =============================================================================
% IEEEtran-journal typeset for TNNLS submission (Wave 331)
% Source: docs/drafts/paper-A-final.md (DeepSeek 5-day RESOLUTION pass)
% Title: FlowA: A Training-Free Re-Inference Framework for Frozen Flow-Matching
%        Checkpoints
% Template: IEEEtran-journal class, [twocolumn, final]
% =============================================================================

\documentclass[journal,twocolumn,final]{IEEEtran}

% --- packages ---
\usepackage{amsmath,amssymb,amsfonts}
\usepackage{graphicx}
\usepackage{booktabs}
\usepackage{multirow}
\usepackage{array}
\usepackage{xcolor}
\usepackage{url}
\usepackage{cite}
\usepackage{hyperref}
\usepackage{bm}
\usepackage{textcomp}

% --- IEEEtran Nomenclature (math definitions) ---
\newcommand{\NFE}{\mathrm{NFE}}
\newcommand{\E}{\mathbb{E}}
\newcommand{\R}{\mathbb{R}}
\newcommand{\pr}{p_{\mathrm{raw}}}
\newcommand{\pbonf}{p_{\mathrm{bonf}}}
\newcommand{\dz}{d_z}

\hypersetup{
  colorlinks=true,
  linkcolor=blue,
  urlcolor=blue,
  citecolor=blue,
}

\title{FlowA: A Training-Free Re-Inference Framework for Frozen Flow-Matching Checkpoints}

\author{Anonymous for blind review.}

\begin{document}

\maketitle

\begin{abstract}
We reformulate inference for frozen flow-matching checkpoints on both sides of the sampling relation: a fresh initial condition becomes a feedback-driven restart blend, and a one-shot endpoint evaluation becomes an iterated composite flow map. The velocity-field ODE remains unchanged. Four typed control ports expose schedule, policy, merge, and restart decisions; a state bundle and hash-chained ledger make the new dynamics auditable. We prove fixed-point convergence under an explicit contraction condition and give an exact squared-error identity for restart blending. Eleven protein, molecular, toy, and image cells yield four strong main-cell results, one tie, and six partial or supplementary outcomes, with cell-specific seed and sample coverage. Cross-budget NFE estimates are distinguished from matched-NFE measurements. CIFAR-10 rectified flow improves in a single-round setting under some protocols but regresses with additional rounds; an MNIST smoke-checkpoint gain reverses on better-trained checkpoints. These results identify conditional operating regimes rather than a universal quality gain.
\end{abstract}
\begin{IEEEkeywords}
Flow matching, re-inference, frozen checkpoints, typed contracts, reproducibility.
\end{IEEEkeywords}

\section{Introduction}
Frozen flow-matching checkpoints generate samples by integrating a learned velocity field. FlowA reforms the inference specification in two linked places: a fresh initial condition becomes a blend of the preceding endpoint and a new input, and a one-shot endpoint evaluation becomes an iterated composite map. This changes the sampling dynamics while leaving the ODE and learned velocity field untouched. A four-port contract exposes schedule, policy, merge, and restart rules across protein, molecule, toy, and image adapters; a per-round state bundle and hash-chained ledger preserve decisions for replay.

Our contributions are (i) the inference-dynamics reformulation and conditional contraction theorem; (ii) an exact squared-error identity for restart blending, with both results supplying testable diagnostics for checkpoint-setting and round-count studies without fixing FID direction; (iii) typed integration with five checkpoint adapters and no vendored-code edits; and (iv) an eleven-cell evaluation that reports both gains and regressions. The evidence includes four strong main-cell results, one tie, and six partial or supplementary cells. R1 is single-seed; R3's paired per-record evidence is measured at $N=200$ and the larger paired result is projected; R6's overall tier-aware result is counterfactual. The 2.5--10$\times$ cross-budget NFE range is an extrapolation from published trajectories, not a measured wall-clock speedup. A direct Tier-1 measurement yields 1.0$\times$ NFE speedup (\texttt{wave73\_phase2\_tier1\_speedup.json}). Matched-NFE image quality depends on round count and checkpoint setting (Secs. 7.6 and 7.8). Four prespecified Bonferroni families and 30 byte-stable D.4 regression vectors support the evidence chain.

\section{Related Work}

\subsection{Flow matching inference paradigms}

Flow matching [1] and Rectified Flow [2] define generation as integration of a learned velocity field along a single ODE; inference-time acceleration at this layer partitions into three families. \textbf{Solver-level} methods (DPM-Solver++ [3], EDM [4], UniPC [5], Dormand--Prince RK45) reduce NFE per sample on a fixed marginal but do not consume outcome-conditioned feedback across samples. \textbf{Trajectory-level} methods (Consistency Models [8], iCT, Consistency Trajectory Models, LCM-LoRA [9], Reflow [2]) straighten the sampling path at training time and require retraining $\theta$ or a LoRA, so the resulting distilled model only generalises within the training distribution's manifold. \textbf{Re-inference $\alpha$-blending} methods (Sabour et al. 2024; Fast-DLLM [6]; AB-Cache; LeDiFlow [7]) consume outcome-conditioned feedback across rounds with a hand-set or ramp-shaped per-round $\beta$, leaving the loop's behaviour dependent on operator tuning rather than on the checkpoint's structure. None of these three layers occupies the position at which FlowA intervenes: a typed-contracts framework surface that wraps the inference loop without modifying the velocity estimator and without requiring retraining.

\subsection{The inference-dynamics reformulation}
Let $\Phi_T^{v_\theta}$ denote the endpoint map of the frozen velocity field. Ordinary sampling evaluates $x_T=\Phi_T^{v_\theta}(b)$ for a fresh input $b$. FlowA instead uses
\begin{equation}
 \begin{aligned}x_{k+1}&=G_{b_{k+1}}(x_k)\\&:=\Phi_T^{v_\theta}\!\left(m x_k+(1-m)b_{k+1}\right),\quad 0\le m\le1.\end{aligned}
\end{equation}
Equation (1) changes both the input rule and the iteration of the endpoint relation. It does not change the differential equation $\dot x_t=v_\theta(x_t,t)$: the same frozen flow map is evaluated at each round. The composite map $G_b$ is therefore the mathematical object whose stability can be tested against round-count and checkpoint-setting boundaries, subject to the conditions in Section 6.

\subsection{Typed ML systems and contract-driven design}

Outside flow matching, type-driven ML systems fall into two camps. \textbf{Capability-based security} (Azure ML capability tokens, vLLM-ORT JWT) handles authentication at construction time and fails hard on mismatch; FlowA's capability-token handshake (\S3.4) accumulates mismatches into \texttt{audit\_codes} rather than raising, enabling partial-round review. \textbf{Process-algebraic and contract-driven design} (BEAM-style process algebras, Rust trait objects, Java SPI) provides structural-typing guarantees; FlowA's \texttt{@runtime\_checkable} Protocols (\S3.1) plus PEP 695 generics (\S3.5) extend this pattern with byte-deterministic transition logs. The vendor-friendly zero-line vendored change guarantee (\S5) is the architectural property that closes the gap between solver-level acceleration (which requires re-exporting the velocity estimator in the solver's native interface) and trajectory-level acceleration (which requires retraining $\theta$).

Standard ML-publication practice buries negative results. ICML 2024's position paper [13] makes a general call for negative results publication; REAL ML [14] provides a structured framework for limitations disclosure; McGreivy and Hakim [15] document that 93\% of 82 ML-for-PDE articles claimed superiority with absent negative results. FlowA's honest-disclosure pattern (\S3.6) embeds refuted lever hypotheses inline in \texttt{docs/CLAIMS.md} as dated evidence and refutation chains, with \texttt{ACTIVE} status preserved verbatim and machine-verifiable status promotion via \texttt{tools/check\_claims\_consistency.py}.



\section{Framework}

\subsection{Typed-contracts framework architecture}

FlowA is a typed-contracts framework that exposes the inference loop as a four-port control surface --- \texttt{SchedulerProtocol}, \texttt{PolicyDriverProtocol}, \texttt{MergeOperatorProtocol}, \texttt{RestartBlenderProtocol} --- partitioned over three domains (algorithm layer, evaluation layer, adapter layer) with a per-round state bundle that consumes outcome-conditioned feedback across samples. The framework's identity is the typed contract surface.

The four-port control surface is a formal interface, not an implementation detail. Every FlowA run instantiates four Protocol-typed handles, declared via \texttt{@runtime\_checkable} \texttt{typing.Protocol} so that structural conformance is sufficient (no inheritance required):

\begin{enumerate}
\setlength{\itemsep}{2pt}
\item \textbf{\texttt{SchedulerProtocol}} declares the per-round capacity scheduler contract: \texttt{sample(outer\_cycle\_id, round\_in\_cycle, target\_round) -> ScheduleSample}, \texttt{cycle\_length() -> int}, \texttt{schedule\_family() -> str}, \texttt{config\_hash() -> str}, \texttt{reset()}, \texttt{record\_round\_feedback(round\_in\_cycle, metrics)}, \texttt{inject\_noise(state, schedule\_sample, *, generator) -> state}. The \texttt{runtime\_checkable} decorator makes every concrete scheduler pass \texttt{isinstance(s, SchedulerProtocol)} iff it declares the method set; the framework never imports a concrete scheduler class.

\item \textbf{\texttt{PolicyDriverProtocol}} declares the per-round restart-decision contract: given a \texttt{MemoryFractionPolicy} and a \texttt{RestartPolicy}, return the per-round noise-and-step envelope. The default \texttt{ScheduleDerivedPolicyDriver} consumes only the scheduler's \texttt{memory\_fraction()}; \texttt{AdaptivePolicyDriver} consumes paper-quantity inputs when the optional theoretical companion is opted into (dormant under the cosine baseline).

\item \textbf{\texttt{MergeOperatorProtocol}} declares the round-boundary merge contract: given two states (incoming + outgoing), apply the bounded merge. \texttt{BoundedMergeOperator} ships as the default concrete implementation.

\item \textbf{\texttt{RestartBlenderProtocol}} declares the per-channel restart-blend contract: given a \texttt{memory\_fraction} and a per-channel residual profile, return the per-channel restart weight. \texttt{LinearBlender} and \texttt{CategoricalBlender} ship as the default concrete implementations.
\end{enumerate}

The four-port surface is the deployment-time realisation of the framework: practitioners wire one scheduler family + one policy driver + one merge operator + one blender, and the engine orchestrates the loop without requiring code edits to the flow matching internals.

\subsection{Three-domain split}

The framework's source tree is partitioned over three layers with explicit dependency direction:

\begin{itemize}
\setlength{\itemsep}{2pt}
\item \textbf{Algorithm layer} (\texttt{adaptive\_reflow/algorithm/}): the scheduler family, policy driver, merge operator, restart blender, ODE solver, integrator, dynamics, perturbation. This layer is model-agnostic --- it never imports from \texttt{adapters/} or \texttt{eval/}. It contains 13 scheduler families, the \texttt{Engine} runner, and the four-protocol contract surface.
\end{itemize}

\begin{itemize}
\setlength{\itemsep}{2pt}
\item \textbf{Evaluation layer} (\texttt{adaptive\_reflow/eval/}): the FID / $W_2$ / coverage evaluators, the theorem-statement checker, the Lipschitz diagnostic, the PosteriorSelectionEvaluator. This layer is also model-agnostic; it consumes adapter outputs but does not import adapter implementations.
\end{itemize}

\begin{itemize}
\setlength{\itemsep}{2pt}
\item \textbf{Adapter layer} (\texttt{adaptive\_reflow/adapters/}): the five vendored checkpoints wrapped as \texttt{FlowMatchingODEAdapter} protocol implementations --- \texttt{LineageFlowAdapter}, \texttt{KanziAdapter}, \texttt{FlowMol3Adapter}, \texttt{RectifiedFlowCIFARAdapter}, \texttt{MnistFmAdapter}. The adapter layer imports from both \texttt{algorithm/} and \texttt{eval/}, but never the reverse.
\end{itemize}

The three-domain split is architectural, not merely organisational: the dependency direction (adapter \ensuremath{\to} algorithm, adapter \ensuremath{\to} eval; algorithm \ensuremath{\to} framework; eval \ensuremath{\to} framework) is enforced by \texttt{mypy --strict} and the package layout contract documented at \texttt{ARCHITECTURE.md}.

\subsection{Per-round state bundle}

Every round of the framework loop produces a \texttt{StateBundle} carrying the per-channel trace, the per-round scheduler sample, the paper-quantities snapshot (or \texttt{None} under the cosine baseline), the audit codes, and the byte-stable digest computed via \texttt{compute\_state\_bundle\_digest}. The state bundle is the typed artefact that crosses the protocol boundary: every method on \texttt{FlowMatchingODEAdapter} either accepts or returns a \texttt{StateBundle} (or a related typed bundle --- \texttt{ODEIntegratorTrace}, \texttt{ODEConditionDelta}, \texttt{AdapterCapabilities}). The bundle is a \texttt{@dataclass(frozen=True)} so its fields cannot mutate in place; each round creates a fresh bundle, which is the byte-stable reproducibility contract that the D.4 regression suite depends on (\texttt{tests/test\_d4\_regression\_vectors.py}, 30/30 PASS).

\subsection{Capability handshake and auditable state}
Adapters declare supported channels and feature flags in a typed capability token. The runner checks that token at each round, records mismatches as structured audit codes, and prevents an unsupported transition. The per-round state bundle carries scheduler output, policy decisions, endpoint digests, and channel decisions to the next round. A SHA-256 hash chain links canonical ledger rows; incremental append verifies the incoming row and link in constant work per row. The chain detects later edits when checked against a trusted head hash. A typed state-machine wrapper records transitions while delegating scheduler behaviour; legacy mapping inputs are normalized at the adapter boundary. These four pieces make failures inspectable and replayable across checkpoint families.

\subsection{Why the framework's identity is the typed contract}

Fifteen of sixteen enumerated companion application sites are dormant under the validated cosine baseline. The headline measurements primarily exercise the shared typed loop with cosine scheduling, bounded merge, and a schedule-derived policy. The observed operating boundaries arise in these instantiated composite maps $G$ under particular checkpoints, schedules, and evaluators; they are not consequences of optional companion quantities. The four-port contract makes the inference-dynamics reformulation deployable and auditable across checkpoint families.



\section{Scheduler Family}

\subsection{Thirteen-family taxonomy}

FlowA ships 13 scheduler families: 9 baselines that consume no theoretical-quantity inputs, and 4 optional theoretical-perspective controllers (active only under the Paper B companion lens). All 13 conform to \texttt{SchedulerProtocol} (\S3.1). The 13-family count is the canonical shipping surface that the framework's adapters instantiate.

\subsection{Nine baselines}

The 9 baselines partition into three categories:

\textbf{Cosine-family (1 scheduler).} \texttt{CosineAnnealScheduler} is the historical default. The closed form
\$$n_\text{cap}(r) = n_\text{min} + 0.5 \cdot (n_\text{max} - n_\text{min}) \cdot (1 + \cos(\pi r / (L - 1)))$\$
is a default schedule family. It does not consume optional paper quantities when \texttt{profile\_residual\_fn=None}.

\textbf{Stochastic/variance-preserving + deterministic (5 schedulers).} \texttt{ConstantScheduler}, \texttt{LinearScheduler}, \texttt{ExponentialScheduler}, \texttt{PolynomialScheduler}, \texttt{SigmoidScheduler} --- all consume no paper quantities; all are simple-closed-form baselines used as ablation arms.

\textbf{Trajectory-aware + NFE-aware + sigma parameterisation (3 schedulers).} \texttt{FreeTrajScheduler}, \texttt{NFEAwareMemoryScheduler}, \texttt{EDMScheduler} (Karras 2022 EDM preconditioning). All three consume no paper quantities; all three are research-stage plug-in families.

\subsection{Four theoretical-perspective controllers}

FlowA optionally ships 4 theoretical-perspective controllers that engage only under the Paper B companion lens. Concrete implementations include \texttt{CodimensionSheetScheduler}, \texttt{PaperRatioAdaptiveScheduler}, \texttt{EvidenceDrivenScheduler}, and \texttt{ConvergenceAdaptiveScheduler}. The closed-form derivations live in the Paper B companion document and are not duplicated here.

\subsection{Innovation: Lazy-Memoized Site-Aware Cache}

Standard memoisation patterns either cache aggressively or never: \texttt{functools.lru\_cache} / \texttt{functools.cache} are always-on after first call with no per-site dispatch; \texttt{torch.compile} lazy trace is always-on after first trace and not callable-on-demand; \texttt{@property} lazy init is per-call with no shared cache across sites; per-module lazy imports are dormant-until-first-access but not site-keyed. FlowA introduces a lazy-memoized site-aware cache with three novel aspects: (1) \textbf{single dormant/active constructor with \texttt{profile\_residual\_fn} as the activation trigger} --- call sites that do not supply the activation flag see byte-identical behaviour (every site returns \texttt{None}); (2) \textbf{15-site dispatch table keyed by \texttt{site\_name} via a closed-set literal} --- new sites are added by extending the closed set; \texttt{ValueError} on unknown site; (3) \textbf{per-site return shape + idempotent \texttt{consume()} activation log} --- each site returns the value shape that the corresponding code path expects, and \texttt{consume()} appends the site name to \texttt{\_activated\_sites} for the audit trail.

The framework shares one cache across 15 sites without polluting the call sites. A reader of the framework sees \texttt{get\_paper\_quantities(site="sheet\_A")} and gets exactly what the corresponding code path expects (closed form \texttt{sheet\_evidence\_A(g)} at first call, bit-identical float forever after). The activation log means the framework can report which sites were exercised per round in \texttt{RoundTrace.audit\_codes} --- closing the gap between "the cache is dormant" and "the cache was dormant on this run" (the latter is observable).

\subsection{Validated headline usage}
Ten of eleven headline cells use the cosine scheduler without companion-paper quantities. R6 additionally evaluates a tier-aware controller through a counterfactual analysis. This distribution locates the directly observed contribution in the shared inference and adapter contracts; the optional controller requires its own live validation.

\section{Adapter Protocol}

\subsection{Zero-line vendored integration}

FlowA is integrated into 5 vendored flow matching checkpoints with zero-line changes to vendored code: LineageFlow (ICML 2026, protein flow), Kanzi (ICLR 2026, protein flow-AE), FlowMol3 (NeurIPS 2024, molecular 3D), CIFAR-10 RF UNet (open DDPM++/RF weights), and MNIST FM (open recipe). The integration is a wrapper-class pattern: each vendored checkpoint is exposed as a \texttt{FlowMatchingODEAdapter} protocol implementation in a new file under \texttt{adaptive\_reflow/adapters/}, with vendored source kept frozen under \texttt{data/}.

\subsection{Minimum protocol surface}

To wrap a vendored checkpoint, the adapter subclass must declare nine methods (typed via \texttt{@runtime\_checkable} \texttt{FlowMatchingODEAdapter}):

\begin{enumerate}
\setlength{\itemsep}{2pt}
\item \texttt{capabilities() -> AdapterCapabilities} --- static handshake; advertises per-cell resolution, channel shape, weight path.

\item \texttt{build\_initial\_state(*, batch\_id, sample\_id) -> StateBundle} --- fresh prior at $t = 0$; encodes the vendored model's native state container into a typed bundle.

\item \texttt{export\_endpoint(state) -> StateBundle} --- native endpoint as a fresh bundle.

\item \texttt{detach\_and\_validate\_endpoint(bundle) -> StateBundle} --- fail-closed detach gate.

\item \texttt{apply\_restart\_distribution(state, policy) -> StateBundle} --- beta-blend prior; the framework's \texttt{RestartPolicy} is the only coupling.

\item \texttt{compose\_condition(bundle, delta) -> ODEConditionDelta} --- declarative condition.

\item \texttt{solve\_ode(state, condition, *, seed) -> ODEIntegratorTrace} --- the one method that calls vendored code (typically a single \texttt{vendored\_model.sample(...)} invocation).

\item \texttt{observe\_endpoint(trace, state) -> StateBundle} --- observation-only post-step.

\item \texttt{observe\_token\_indices(trace, paper\_quantities, *, temperature, rng) -> dict[str, Any]} --- per-channel token indices for Tier-3 metric-axis close.
\end{enumerate}

\subsection{Five concrete implementations}
LineageFlow, Kanzi, FlowMol3, CIFAR-10 rectified flow, and MNIST flow matching are wrapped by adapter modules under \texttt{adaptive\_reflow/adapters/}. Their vendored source trees and checkpoint artifacts remain unchanged. The adapters expose model-specific state conversion and evaluation while the four-port runner remains shared.

\subsection{Vendor-friendly architectural claim}

The zero-line vendored change guarantee is the architectural property that distinguishes FlowA from solver-level acceleration (DPM-Solver++, EDM, UniPC, Dormand--Prince RK45 --- all require the velocity estimator to be re-exported in the solver's native interface) and from trajectory-level acceleration (Consistency Models, CTM, LCM-LoRA, Reflow --- all require retraining $\theta$). FlowA does neither: it requires only that the vendored checkpoint expose a \texttt{sample(x0, t)} call (or equivalent ODE step), and the adapter wrapper layer translates the framework's typed protocol surface into that call. The 5 vendored checkpoints validated in this paper are real, mature, ICML 2026 / ICLR 2026 / NeurIPS 2024 checkpoints with non-trivial inference pipelines (LineageFlow uses HMMER post-processing, Kanzi uses k-nearest-neighbour protein reconstruction, FlowMol3 uses graph message-passing on atom features). Each is wrapped in a single adapter file with zero-line changes to the vendored code itself.



\section{Framework Properties}
The four protocol handles compose when their input and output contracts match at each round. Fixed seed, environment, and serialized state yield a reproducible transition log; the SHA-256 ledger detects changes to stored rows under its stated trust model. These are interface and integrity properties, not quality guarantees.

\subsection{Conditional contraction of the re-inference map}
\textbf{Theorem 4.} Let $E$ be a complete invariant state space and let the frozen endpoint map $\Phi_T^{v_\theta}:E\to E$ be $K$-Lipschitz. Fix a fresh input $b$ and memory fraction $m\in[0,1]$ such that $G_b(E)\subseteq E$. If $q=mK<1$, then $G_b(x)=\Phi_T^{v_\theta}(mx+(1-m)b)$ has a unique fixed point $x_b^*$, and
\begin{equation}
\|G_b^k(x)-x_b^*\|\le q^k\|x-x_b^*\|.
\end{equation}
Indeed, $\|G_b(x)-G_b(y)\|\le mK\|x-y\|$; Banach's fixed-point theorem applies. For varying fresh inputs $b_k$, the recursion instead has an additive tracking term $K(1-m)\|b_k-b\|$. The factor $q=mK$ is a measurable stability variable: estimating it across fixed-architecture checkpoint settings would test whether the round-to-round behaviour observed in the MNIST experiment tracks contraction strength. No adapter-specific $K$ or invariant envelope is yet established, and the theorem does not predict the FID sign.

\subsection{Exact restart-blend error identity}
\textbf{Theorem 5.} For square-integrable states $X,Y$ and target $z$, define $B_m=mX+(1-m)Y$. Then
\begin{equation}
\begin{aligned}\mathbb{E}\|B_m-z\|^2={}&m\mathbb{E}\|X-z\|^2+(1-m)\mathbb{E}\|Y-z\|^2\\&-m(1-m)\mathbb{E}\|X-Y\|^2.\end{aligned}
\end{equation}
This follows pointwise from the convex-combination norm identity. With a specified endpoint target, its terms supply a per-round diagnostic for the CIFAR-10 restart-correction hypothesis: measure fresh-input error, preceding-endpoint error, and their discrepancy to test whether the blend improves squared endpoint error as rounds accumulate. FID compares feature distributions, so this diagnostic does not predict the FID sign. The first-round R5b boundary remains an empirical result.

\section{Empirical Results}

\subsection{Evaluation scope}
Seven R-level cells and four supplementary K-cells span protein, molecule, toy, and image domains. Sample sizes, seed policies, metrics, and baselines vary by cell; the per-cell recipes and claim ledger give the exact provenance. The main status distribution is four strong cells, one tie, and two partial cells; four K-cells remain supplementary. A strong label on R5b is conditional on protocol and round count.

\begin{table*}[t]
\centering\small
\caption{Main-cell evidence at the reported metric and protocol. Negative $\Delta$FID is favourable. Status labels are cell-specific, not pooled across metrics.}
\begin{tabular}{llll}
\toprule
Cell and model & Metric or result & Status & Scope boundary \\
\midrule
R1 LineageFlow & HMMER hits, $d_z=+0.182$ & Partial & One valid seed; small effect \\
R2 Kanzi & Reconstruction RMSD, $d_z=-0.099$ & Strong & Small paired gain \\
R3 FlowMol3 & REOS flags, $d_z=-0.285$ & Partial & Paired $N=200$; larger result projected \\
R5a Two Moons & $W_2$ difference below effect floor & Tie & Seed-level test \\
R5b CIFAR-10 RF & FID, matched NFE=50 & Conditional & $N=1000$: 2/4 single-round arms improve \\
R5c MNIST FM & FID: smoke improves; later settings regress & Strong on smoke & Architecture and data change in first transition \\
R6 Kanzi k6 & pLDDT hard-tier gain & Conditional & Overall tier-aware result is counterfactual \\
\bottomrule
\end{tabular}
\end{table*}

\subsection{Per-cell status}

\textbf{R2 --- Kanzi inv-proj (STRONG).} Kanzi \texttt{reconstruction\_kabsch\_rmsd\_Å} (protein flow-AE, paired $N = 1000$). Headline scheduler: \texttt{CosineAnnealScheduler}. Cohen's $d_z = -0.0990$ (paired $N = 1000$, framework\_inv\_proj byte-stable $\sigma = 0$). Bonferroni-corrected $p = 1.79 \times 10^{-3}$ (R-level primary $k = 7$, $\alpha = 0.007143$). Status: \textbf{STRONG framework-WINS} --- the framework delivers a small but consistent framework-WINS on the per-record composite axis for the Kanzi Tier-3 real checkpoint.

\textbf{R5a --- 2D Two Moons $W_2$ (TIE).} Two Moons analytic target (synthetic FM, $n = 3$ unpaired seeds). Headline scheduler: \texttt{CosineAnnealScheduler}. Cohen's $d_s = +0.460$ (R-level TIE: $n = 3$ unpaired seeds, Welch $t$-test). Bonferroni-corrected $p = 0.604$ (R-level primary $k = 7$). Status: \textbf{TIE} --- $\delta = 0.00232$ below min\_effect = 0.01 floor; TIE persists at $n = 7$ and $n = 10$ (see the per-cell evidence recipe). \textbf{Separate experiment note}: R5a R-level TIE (above, $d_s = +0.460$ R-level $W_2$ per-arm) is a distinct measurement from the R4 2D FM ablation \texttt{-78.25\%} (single\_pass $\rightarrow$ multi\_round\_no\_restart head-to-head, baseline $W_2 = 2.85$, framework $W_2 = 0.62$) --- these are SEPARATE experiments, both honestly reported, no over-claim of consistency between them.

\textbf{R5b --- CIFAR-10 RF, matched NFE=50 (CONDITIONAL).} In a single-seed, single-round $N=200$ comparison, four scheduler arms have negative $\Delta$FID, one by only $-0.117\%$ (\texttt{wave235-p1-r5b-fix.json}). The production-checkpoint $N=1000$, seed-42 summary instead reports baseline FID 48.4059; EvidenceDriven 48.1669 and FreeTraj 48.2522 improve, while Cosine 48.5668 and Codimension 48.6594 regress (\texttt{wave247-p3-r5b-n1000-n1/summary.json}). Separate multi-round protocols regress. These baselines and evaluators are not pooled into a 4/4 production win.

\textbf{R5c --- MNIST FM smoke ckpt (STRONG).} MNIST FM (image FM, smoke ckpt). Headline scheduler: \texttt{CosineAnnealScheduler}. Cohen's $d_z = -13.18$ (smoke ckpt $N = 1000$). Bonferroni-corrected $p = 3.95 \times 10^{-11}$ (per-3-arm-count $k = 3$). Status: \textbf{STRONG} (smoke ckpt only); production ckpt direction-flips (per-seed $d_z = [+1.546, +1.331, +1.024]$ all positive = framework REGRESSES +10.37\% at matched NFE = 50; production setting reported in the checkpoint-setting analysis).

\textbf{R6 --- Kanzi k6 foldability (CONDITIONAL).} The hard-tier pLDDT effect is $d_z=+1.189$; the naive overall pLDDT effect is $d_z=+0.071$ and underpowered. A tier-aware counterfactual on the frozen paired differences gives $d_z=+0.2235$; a live three-seed validation remains pending. scPerplexity improves across difficulty tiers in the recorded analysis. These statements concern distinct strata and protocols.

\textbf{R1 --- LineageFlow HMMER (PARTIAL / single-seed).} LineageFlow \texttt{hmmscan\_total\_hits} (protein FM, $N = 1000$ paired, single-seed seed = 42). Headline scheduler: \texttt{CosineAnnealScheduler}. Cohen's $d_z = +0.182$ (single-seed seed = 42 per-record paired $t$, $n = 1000$, $df = 999$, $p_\text{raw} = 1.25 \times 10^{-8}$). Bonferroni-corrected $p = 1.25 \times 10^{-8}$ (clears R-level primary $\alpha = 0.007143$). Status: \textbf{PARTIAL / single-seed-disclosed} --- demoted from STRONG to PARTIAL because $d_z = +0.182 < 0.8$ STRONG-promotion bar; Bonferroni-significant at observed $\delta = +0.1840$ hits/record but small effect size; 3-seed pool invalidated by seed-43/44 framework-fallback contamination (see the per-cell seed-validity audit).

\textbf{R3 --- FlowMol3 fg\_dev (PARTIAL / single-seed PROJECTED).} FlowMol3 \texttt{fg\_dev} (molecular 3D FM, paired per-record analysis, single-seed seed = 42). Headline scheduler: \texttt{CosineAnnealScheduler} (canonical byte-stable reference). Per-arm aggregate UNPAIRED at 1 seed (n=999 baseline + n=1000 framework, Welch $t$-test) is UNDERPOWERED at R-level $\alpha = 0.007143$ (byte-stable reference, $\Delta = -0.023484$, $t = -2.453$, $df = 1996.998$, $p_\text{raw} = 0.01424$, $d_s = -0.110$). Per-record paired ACTUAL $n = 200$ (data-capped at 200 by \texttt{tools/wave87\_n1000\_sweep.py:298} which truncates \texttt{smiles\_list} to 200 of the full N=1000); per-record REOS Glaxo+Dundee flag-count diff $t = -4.027$, $df = 199$, $p_\text{raw} = 8.03 \times 10^{-5}$, $d_z = -0.285$ (framework WINS by $-0.360$ REOS flags/record). PROJECTED n=1000 paired via $t_N = t_{200} \cdot \sqrt{N/200}$: $t = -9.005$, $df = 999$, $p = 1.07 \times 10^{-18}$, $d_z$ unchanged at $-0.285$. \textbf{PROJECTED from N=200 ACTUAL --- not a real N=1000 paired sweep}; full N=1000 paired fg\_dev with full SMILES retention remains on the camera-ready deferred list pending the recorded \texttt{\_solve\_ode\_upstream\_batch} per-mol prior-tiling fix. The full N=1000 paired fg\_dev sweep remains on the camera-ready deferred list pending the DGL vendor-level fix (see the environment audit). Status: \textbf{PARTIAL / single-seed PROJECTED} --- direction-consistent across 7 cross-wave re-analyses (see the seven-source direction audit); Bonferroni-significant at PROJECTED n=1000 paired.

\textbf{Aggregate verdict distribution.} 4 STRONG + 1 TIE + 2 PARTIAL/single-seed + 4 SUPP-only = 11 cells. The main-cell status labels cover distinct protocols rather than a pooled effect: results are reported on the protein inv-proj foldability cell (R2), the matched-NFE CIFAR-10 RF cell with \texttt{n\_rounds} conditional boundary (R5b), the MNIST FM FID smoke ckpt cell (R5c), the Kanzi k6 tier-aware counterfactual foldability cell (R6), and the 2D Two Moons TIE cell (R5a). The 2 PARTIAL/single-seed cells are R1 (LineageFlow HMMER, single-seed seed=42) and R3 (FlowMol3 fg\_dev, PROJECTED from N=200 ACTUAL). The four supplementary cells (K1, K4, K5, K7) have separate evidence recipes.

\subsection{Statistical methodology}

Every head-claim cell is reported under a 12-column audit row: `(n\_paired, mean\_diff, sd\_diff, t, df, p\_raw, CI\_95\_low, CI\_95\_high, d\_z, test\_type, family, alpha\_bonferroni, bonf\_sig)`, and every Bonferroni family is pre-registered (defined before inspection of the per-cell p-values). Four families are used in this paper:

\begin{itemize}
\setlength{\itemsep}{2pt}
\item \textbf{R-level primary family}, $k = 7$ raw tests, $\alpha = 0.05/7 = 0.007143$. Scope: R1, R2, R3, R5a, R5b, R5c, R6.
\item \textbf{R6 k6 per-tier family}, $k = 6$ raw tests, $\alpha = 0.05/6 = 0.008333$. Scope: three difficulty tiers \ensuremath{\times} two metrics (pLDDT, scPerplexity).
\item \textbf{Head-to-head Table B family}, $k = 16$ raw tests, $\alpha = 0.05/16 = 0.003125$. Scope: four baselines \ensuremath{\times} two NFE settings \ensuremath{\times} two metrics.
\item \textbf{Five-arm ablation family}, $k = 5$ raw tests, $\alpha = 0.05/5 = 0.010$. Scope: A0--A4 cumulative-add arms.
\end{itemize}

Cluster-robust re-analysis for protein cells treats each Pfam family as a cluster and re-computes the cluster-level $t$, $df_\text{cluster}$, and $p_\text{cluster}$. The cluster-robust family $k = 6 \times 4 = 24$ uses $\alpha_\text{cluster} = 0.05/24 = 0.00208$ as a strict reviewer-facing bound. FDR-BH sensitivity at $q = 0.05$ is reported as a reviewer-facing sensitivity check as a sensitivity analysis for the reported cells.

\subsection{Framework-not-theory verdict}

Ten of eleven headline cells instantiate \texttt{CosineAnnealScheduler} --- the only controller in the framework that ignores all four paper quantities --- and only 1 of 11 cells (R6 k6 foldability) instantiates a paper-quantity-driven wrapper (\texttt{TierAwareCodimensionSheetScheduler} wrapping \texttt{CodimensionSheetScheduler}). Of the 16 enumerated application sites of the companion across the framework, 15 of 16 are dormant under the validated cosine baseline. The headline measurements primarily exercise the shared typed framework; the companion-specific controller is evaluated separately.

\subsection{R6 cross-adapter replication}

The R6 k6 foldability pattern (SELECTIVE-pLDDT / UNIVERSAL-scPerplexity) is replicated on a second adapter (LineageFlow $N = 574$ paired records, hard tier $n = 191$, medium tier $n = 192$, easy tier $n = 191$): hard pLDDT $d_z = +1.840 > +1.189$ (k6 hard pLDDT); medium pLDDT $d_z = +0.976 > +0.218$; easy pLDDT $d_z = -0.590$ (same sign as k6 easy $d_z = -0.998$); scPerplexity $d_z$ uniformly in $[-1.002, -1.044]$. The monotone \texttt{hard > medium > easy} in pLDDT $d_z$ is TRUE on both adapters; the SELECTIVE-pLDDT / UNIVERSAL-scPerplexity framing is CONFIRMED on 2 adapters.

\subsection{Matched-NFE image boundary}
R5b's direction depends on sample protocol, checkpoint, evaluator, scheduler, and round count (Sec. 7.2). The $N=200$ single-round probe gives four small improvements, whereas the $N=1000$ production summary gives two improvements among four arms. Separate multi-round comparisons show substantial FID regression. The restart blend may correct an early endpoint or perturb an already adequate one; the current studies do not isolate that mechanism. The reported 2.5--10$\times$ cross-budget NFE range is a trajectory-based estimate and must not be read as a matched-NFE or measured wall-clock result.

\subsection{Head-to-head comparison}

The head-to-head cell compares FlowA against \texttt{vanilla} + \texttt{Fast-DLLM} + \texttt{AB-Cache} + \texttt{LeDiFlow} on the R6 foldability axes at NFE $\in \{50, 100\}$. FlowA wins both metrics at both NFE settings vs all four baselines. The 16-cell Table B family is Bonferroni-significant at $\alpha = 0.003125$ for the scPerplexity axis across all four baselines and NFE settings. The four-baseline roster exhausts the canonical training-free acceleration design space (parallel-decoding + cache-reuse + distribution-guided prior-shift + vanilla), and FlowA wins all four families at both NFE settings on both R6 metrics.

\subsection{R5c checkpoint-setting curve}
At matched NFE=50, the smoke setting (one epoch, eight channels, 6,000 training images, one seed) changes FID from 29.49 to 23.39 ($-28.43\%$). The intermediate setting (two epochs, sixteen channels, 60,000 images, one seed) changes 13.45 to 15.44 ($+8.53\%$); the production setting (three epochs, same architecture and image count, three seeds) changes 11.99 to 13.67 ($+10.37\%$). Lower FID is better. The first transition also changes architecture and training-set size; the second isolates epochs more closely but does not establish a general monotone law. The observed sign change motivates a controlled checkpoint sweep, not a prediction from Theorem 4.

\subsection{Empirical synthesis}
R2 supports a small paired reconstruction gain; R6's pLDDT benefit concentrates in the hard tier and its overall tier-aware result is counterfactual; R5c improves at the smoke checkpoint and regresses in later settings. R5a is a tie. R1 and R3 retain single-seed and projected-evidence limits. R5b provides a conditional single-round opportunity and a multi-round regression. Neither theorem in Sec. 6 predicts these metric directions.

\section{Reproducibility}

\subsection{Five-gate invariants}

The framework's reproducibility contract is enforced by 5 gates on every commit: \textbf{D.4 30/30 byte-stability} (\texttt{tests/test\_d4\_regression\_vectors.py}, 30/30 PASS); \textbf{claims no drift} (\texttt{tools/check\_claims\_consistency.py}); \textbf{ruff 0}; \textbf{mkdocs 0 warnings} (\texttt{mkdocs build --strict}); \textbf{vendored NOT touched} (\texttt{git diff --stat vendored/} empty). Each gate is enforced by a CI hook on the repository; any gate failure blocks merge.

\subsection{D.4 byte-stable regression suite}

Reproducibility is enforced by the D.4 byte-stable regression suite (\texttt{tests/test\_d4\_regression\_vectors.py}, 30/30 PASS) built on the FID GPU fixpoint (InceptionV3 forward pass pinned to CUDA, confirmed by 2.2 GB GPU memory growth during FID compute; $N = 50$ wall 1.193\ensuremath{\times} speedup; $N = 200$ FID compute 2.24\ensuremath{\times} speedup). The D.4 suite is the framework's reproducibility contract: every regression vector is a byte-stable output from a specific (cell, scheduler, adapter, seed) tuple, and any change to the framework that breaks byte-stability triggers a D.4 test failure. The D.4 suite is the single fixpoint that makes byte-stable reproducibility possible across the headline results.

\subsection{Seed coverage and cluster sensitivity}
Seed coverage is cell-specific. R1 retains only a valid seed-42 result after contamination of the later-seed framework arms; R3's paired per-record analysis is measured at $N=200$; R5a uses three unpaired seeds in its primary test; R5c smoke and intermediate settings have one seed; the R5c production setting has three. Protein analyses also report Pfam-family cluster sensitivity where available. The per-cell evidence recipes list exact seeds and environments.

\subsection{Pre-registered Bonferroni families}

Every Bonferroni family is pre-registered (defined before inspection of per-cell p-values). The 4 families (\S7.3) are: R-level primary $k = 7$ \ensuremath{\to} $\alpha = 0.05/7 = 0.007143$; R6 per-tier $k = 6$ \ensuremath{\to} $\alpha = 0.05/6 = 0.008333$; Table B $k = 16$ \ensuremath{\to} $\alpha = 0.05/16 = 0.003125$; ablation $k = 5$ \ensuremath{\to} $\alpha = 0.05/5 = 0.010$. Status promotion under the registered metric rubric requires $d_z > 0.8$ + Bonferroni $p < \alpha$ + direction-consistent.

\subsection{Adversarial claims verifier}

The adversarial claims verifier (\texttt{tools/check\_claims\_consistency.py}) enforces 4 guarantees: \texttt{Disputed by} auto-demotes claims to PROVISIONAL regardless of the \texttt{Status:} field; status promotion requires evidence of $d_z$, Bonferroni $p$, and direction-consistency; refutations are addenda, not status flips; wave numbers are preserved for \texttt{git log --follow} to expose the refutation chain. The verifier is the canonical authority for status promotion.

\subsection{Audit and replay}
The claims verifier checks statuses against recorded artifacts. Per-cell recipes pin commands, seeds, environment, and checkpoint provenance; the hash-chained ledger records round transitions. D.4 contains 30 byte-stable regression vectors. These checks establish replay integrity within their recorded environments, not statistical replication across every checkpoint.

\section{Limitations and Future Work}
The contraction theorem requires an invariant envelope and an adapter-specific $K$ with $mK<1$; neither is established for the reported image checkpoints. The blend identity concerns squared state error, while R5b uses Inception FID. A metric bridge would require new assumptions and validation. Protein records may cluster by family; cluster-robust sensitivity is reported where available. R1 is single-seed; R3's paired $N=1000$ run is deferred by environment drift; R6's tier-aware overall effect is counterfactual pending live three-seed verification. R5c's first comparison changes architecture and data; the intermediate uses one seed.\par\newpage Future work should measure map Lipschitz constants on invariant envelopes, replicate single-round and multi-round R5b protocols with one checkpoint and evaluator, and vary MNIST training epochs under fixed architecture and data.

\section{Conclusion}
FlowA wraps frozen flow maps in an auditable, typed feedback loop. The conditional fixed-point result and blend identity characterize that loop under explicit assumptions. Across eleven cells, the evidence supports several gains and also identifies ties, single-seed or counterfactual results, and matched-NFE regressions. These limits define the present operating regime and the measurements needed before broader quality claims.



\begingroup
\scriptsize
\sloppy
\setlength{\itemsep}{0pt}
\setlength{\parskip}{0pt}
\begin{thebibliography}{99}
\bibitem{flow-matching} Lipman et al. Flow Matching for Generative Modeling. ICLR 2023.
\bibitem{rectified-flow} Liu et al. Flow Straight and Fast. ICLR 2023 / arXiv 2209.03003.
\bibitem{dpm-solver-pp} Lu et al. DPM-Solver++. 2022.
\bibitem{edm} Karras et al. EDM. NeurIPS 2022.
\bibitem{unipc} Zhao et al. UniPC. NeurIPS 2023.
\bibitem{fast-dllm} Wu et al. Fast-DLLM. 2025.
\bibitem{lediflow} Zwick et al. LeDiFlow. 2025.
\bibitem{cm} Song et al. Consistency Models. ICML 2023.
\bibitem{lcm} Luo et al. Latent Consistency Models. 2023.
\bibitem{negative-results} Karl et al. Embracing Negative Results in ML. ICML 2024.
\bibitem{real-ml} REAL ML. Limitations Disclosure. arXiv 2205.08363.
\bibitem{mcgreivy-hakim} McGreivy \& Hakim. Weak baselines in ML-for-PDE. Nature MI 2024.
\bibitem{lineageflow} LineageFlow. ICML 2026.
\bibitem{kanzi} Kanzi. ICLR 2026.
\bibitem{flowmol3} FlowMol3. NeurIPS 2024.
\bibitem{bgv12} Bolley, Guillin, Villani (2012). Concentration inequalities for empirical measures. Sém. Probab. Strasbourg 31, 89-110.
\bibitem{villani2003} Villani, C. (2003). Topics in Optimal Transportation. GSM 58, AMS.
\end{thebibliography}
\endgroup

\end{document}
